\section{Introduction}

Graph neural networks (GNNs) are now a standard approach for learning from molecular and protein graphs because they encode local neighborhood structure together with higher-order relational context \cite{kipf2016semi,gilmer2017neural,hamilton2017inductive,wu2020comprehensive,zhou2020graph}. In biomolecular graph classification, this combination matters because predictive evidence may depend simultaneously on local biochemical motifs and broader graph-level organization. Benchmark suites such as PROTEINS, DD, ENZYMES, MUTAG, AIDS, and Mutagenicity therefore remain useful as controlled testbeds for methodology-oriented graph classification studies rather than as direct proxies for downstream biological discovery.

The main methodological difficulty is that deeper message passing is not automatically beneficial in these settings. Repeated propagation can suppress informative intermediate states through over-smoothing and related depth effects \cite{li2018deeper,oono2020simple,chen2023revisiting}, while over-squashing can cause information from distant nodes to be compressed into uninformative bottlenecks \cite{alon2020bottleneck,topping2021understanding}. Residual connections, normalization, and regularization alleviate some of this degradation \cite{he2016deep,bresson2017residual,zhao2020pair,cai2021graphnorm,rong2019dropedge,chenSimpleDeepGraph2020}, but most residual designs still reuse information along the same branch. As a result, they improve optimization without directly testing whether alternative information-flow topologies preserve more useful historical signals.

Related graph-learning work has already explored skip connections, dense aggregation, APPNP-style propagation, Jumping Knowledge readout, and hierarchical pooling \cite{xu2018representation,klicpera2018combining,gao2019graph,ying2018hierarchical}. These methods preserve information from multiple depths or across pooling hierarchies, yet they do not directly isolate whether exchanging hidden states across branches or reinjecting pooled graph summaries across block boundaries is preferable to ordinary same-path residual reuse. No prior study directly compares these reuse topologies under a controlled protocol---this is the comparison the present work provides.

This paper studies Cross-Residual Graph Neural Networks (CR-GNN), a controlled model family that compares plain stacking, node-level residual reuse, node-level cross-residual exchange, graph-level residual propagation, and graph-level cross-summary exchange under shared graph-classification backbones. The study is organized around three specific research questions:

\begin{enumerate}
    \item \textbf{When does cross-residual interaction help?} Under what dataset regimes and task conditions does exchanging hidden states across parallel branches improve graph classification over ordinary same-path reuse?
    \item \textbf{When does ordinary residual reuse remain the stronger default?} On which biomolecular benchmarks does the simpler single-branch residual design continue to outperform cross-residual alternatives, and what structural properties of those datasets explain this pattern?
    \item \textbf{How does residual routing design change the answer?} Beyond the binary question of whether an extra reuse path exists, how do different filtering strategies---learnable dense routing, top-k selection, and sparse suppression---affect the relative performance of residual and cross-residual architectures?
\end{enumerate}

Framing the study around these three questions is important for the paper's methodological scope. Rather than claiming that one additional pathway always yields a better GNN, the goal is to reveal \emph{when and how} residual routing helps---identifying which reuse topology preserves discriminative historical structural signal reliably once backbone, data split, and training protocol are controlled.

The work is organized as a biomolecular graph-classification methods study. The main evidence comes from a protein-oriented benchmark core built around PROTEINS, DD, and ENZYMES, together with supplementary molecular robustness datasets. This design keeps the paper close to biologically meaningful graph settings while maintaining a reproducible and controlled benchmark scope.

\subsection{Contributions}

The main contributions of this work are:

\begin{itemize}
    \item \textbf{A controlled comparison framework for reuse topologies.} Five information-flow designs---plain stacking, node-level residual reuse, node-level cross-residual exchange, graph-level residual propagation, and graph-level cross-summary exchange---are compared under identical backbones, data splits, and training protocols. By holding all other factors constant, the framework isolates the effect of reuse topology on performance, consistency, and representational dynamics.

    \item \textbf{The discovery of systematic trade-offs between peak accuracy, cross-dataset consistency, and representational continuity.} No single topology dominates all three dimensions. Graph-level residual propagation achieves the highest peak accuracy but with the greatest rank volatility; cross-residual architectures are the most consistent performers; and the strongest representational continuity does not predict the best classification accuracy. This three-way tension, quantified through winner-count and rank-stability metrics, provides a new lens for evaluating GNN architectures beyond single-number leaderboard comparisons.

    \item \textbf{A depth-scaling mechanism analysis revealing a CKA--performance paradox.} Layer-wise CKA, cosine similarity, and gradient norm tracking across depths shows that the architecture with the weakest representational continuity achieves the most frequent wins---challenging the assumption that stronger layer-wise stability predicts better performance. This finding suggests that structured representational change, when coupled with appropriate context reinjection, can be productive rather than destructive.

    \item \textbf{Evidence that routing strategy is a meaningful secondary design axis.} Sparse routing combined with cross-residual architectures forms the most frequently optimal configuration across representative targets, and routing effects are target-dependent. This reframes the discussion from binary architectural choices to the joint selection of topology and routing strategy.
\end{itemize}
